\documentclass[11pt, letterpaper]{article}

\usepackage[utf8]{inputenc}
\usepackage[spanish]{babel}
% --- Paquetes de Diseño y Color ---
\usepackage[margin=2.5cm]{geometry}
\usepackage[svgnames]{xcolor} 

% --- Configuración del Fondo (Método Alternativo) ---
\definecolor{fondoGris}{HTML}{F0FFF4} 
% Aplicamos el color al lienzo completo
\AtBeginDocument{\pagecolor{fondoGris}}

% --- Tipografía ---
\usepackage{charter}
\usepackage{parskip}

% --- Información ---
\title{\textbf{HMF Repaso}}
\author{Darío Sheward Castillo}
\date{\today}

\begin{document}

\maketitle

\section*{Objetivo}
Estudiar sistemas de partículas con interaciones a distacia a través del modelo HMF (Hamiltonian Mean Field). Trataremos con partículas sometidas a un potencial simple y periódico. 

\subsection*{Visualización del problema}
En cuanto tratamos de visualizar la problemática a trabajar podemos vernos tentados a relacionar el ángulo descrito como la posición de la partícula en una circunferencia. Sin embargo, esta interpretacioón resulta útil especificamente para visualizar magnetización del sistema en un instante de tiempo. Mientras que la el visualizar a las partículas en un espacio unidimensional donde poseen spin $\theta_i$ conserva la idea de interacción de largo alcance. 


\end{document}